% Figure: Pipeline-aligned diagnostic objects.
% Single-column (figure). Include with: % Figure: Pipeline-aligned diagnostic objects.
% Single-column (figure). Include with: % Figure: Pipeline-aligned diagnostic objects.
% Single-column (figure). Include with: % Figure: Pipeline-aligned diagnostic objects.
% Single-column (figure). Include with: \input{figures/decomposition}

\begin{table}[t]
\centering
\small
\setlength{\tabcolsep}{5pt}
\renewcommand{\arraystretch}{1.35}
\begin{tabular}{@{}p{0.19\linewidth}p{0.27\linewidth}p{0.42\linewidth}@{}}
\toprule
\textbf{Observed object} & \textbf{Signal family} & \textbf{Diagnostic question} \\
\midrule
Answer samples
&
Answer-state uncertainty
&
Do repeated generations disagree, or has the answer surface become
artificially stable? \\
\addlinespace[0.25em]
Retrieved passages
&
Evidence-state uncertainty
&
Does the retrieved text locally entail, contradict, or fail to support the
generated answer? \\
\addlinespace[0.25em]
Graph trace
&
Retrieval-state uncertainty
&
Which anchors, typed triples, and paths made the answer reachable, and where
does the graph abstain? \\
\bottomrule
\end{tabular}

\vspace{0.45em}
\begin{minipage}{0.96\linewidth}
\footnotesize
\textit{The families are not independent causes of correctness.}
They are separately logged views of the same KG-RAG run: generated text,
retrieved evidence, and symbolic retrieval state.
\end{minipage}

\caption{%
  \textbf{Uncertainty is treated as an audit problem over pipeline objects.}
  Answer-state metrics see only sampled answers.
  Evidence-state metrics compare the answer with retrieved passages.
  Retrieval-state diagnostics inspect the graph trace itself: anchors,
  subject--relation--object triples, paths, and abstentions.
  This is why the families are reported separately.
  This table defines the objects; Table~\ref{tab:failure_modes} states how
  each family behaves across retrieval regimes, and
  Table~\ref{tab:failure_taxonomy} maps their joint outcomes into a
  $2{\times}2$ failure taxonomy.%
}
\label{tab:decomposition}
\end{table}


\begin{table}[t]
\centering
\small
\setlength{\tabcolsep}{5pt}
\renewcommand{\arraystretch}{1.35}
\begin{tabular}{@{}p{0.19\linewidth}p{0.27\linewidth}p{0.42\linewidth}@{}}
\toprule
\textbf{Observed object} & \textbf{Signal family} & \textbf{Diagnostic question} \\
\midrule
Answer samples
&
Answer-state uncertainty
&
Do repeated generations disagree, or has the answer surface become
artificially stable? \\
\addlinespace[0.25em]
Retrieved passages
&
Evidence-state uncertainty
&
Does the retrieved text locally entail, contradict, or fail to support the
generated answer? \\
\addlinespace[0.25em]
Graph trace
&
Retrieval-state uncertainty
&
Which anchors, typed triples, and paths made the answer reachable, and where
does the graph abstain? \\
\bottomrule
\end{tabular}

\vspace{0.45em}
\begin{minipage}{0.96\linewidth}
\footnotesize
\textit{The families are not independent causes of correctness.}
They are separately logged views of the same KG-RAG run: generated text,
retrieved evidence, and symbolic retrieval state.
\end{minipage}

\caption{%
  \textbf{Uncertainty is treated as an audit problem over pipeline objects.}
  Answer-state metrics see only sampled answers.
  Evidence-state metrics compare the answer with retrieved passages.
  Retrieval-state diagnostics inspect the graph trace itself: anchors,
  subject--relation--object triples, paths, and abstentions.
  This is why the families are reported separately.
  This table defines the objects; Table~\ref{tab:failure_modes} states how
  each family behaves across retrieval regimes, and
  Table~\ref{tab:failure_taxonomy} maps their joint outcomes into a
  $2{\times}2$ failure taxonomy.%
}
\label{tab:decomposition}
\end{table}


\begin{table}[t]
\centering
\small
\setlength{\tabcolsep}{5pt}
\renewcommand{\arraystretch}{1.35}
\begin{tabular}{@{}p{0.19\linewidth}p{0.27\linewidth}p{0.42\linewidth}@{}}
\toprule
\textbf{Observed object} & \textbf{Signal family} & \textbf{Diagnostic question} \\
\midrule
Answer samples
&
Answer-state uncertainty
&
Do repeated generations disagree, or has the answer surface become
artificially stable? \\
\addlinespace[0.25em]
Retrieved passages
&
Evidence-state uncertainty
&
Does the retrieved text locally entail, contradict, or fail to support the
generated answer? \\
\addlinespace[0.25em]
Graph trace
&
Retrieval-state uncertainty
&
Which anchors, typed triples, and paths made the answer reachable, and where
does the graph abstain? \\
\bottomrule
\end{tabular}

\vspace{0.45em}
\begin{minipage}{0.96\linewidth}
\footnotesize
\textit{The families are not independent causes of correctness.}
They are separately logged views of the same KG-RAG run: generated text,
retrieved evidence, and symbolic retrieval state.
\end{minipage}

\caption{%
  \textbf{Uncertainty is treated as an audit problem over pipeline objects.}
  Answer-state metrics see only sampled answers.
  Evidence-state metrics compare the answer with retrieved passages.
  Retrieval-state diagnostics inspect the graph trace itself: anchors,
  subject--relation--object triples, paths, and abstentions.
  This is why the families are reported separately.
  This table defines the objects; Table~\ref{tab:failure_modes} states how
  each family behaves across retrieval regimes, and
  Table~\ref{tab:failure_taxonomy} maps their joint outcomes into a
  $2{\times}2$ failure taxonomy.%
}
\label{tab:decomposition}
\end{table}


\begin{table}[t]
\centering
\small
\setlength{\tabcolsep}{5pt}
\renewcommand{\arraystretch}{1.35}
\begin{tabular}{@{}p{0.19\linewidth}p{0.27\linewidth}p{0.42\linewidth}@{}}
\toprule
\textbf{Observed object} & \textbf{Signal family} & \textbf{Diagnostic question} \\
\midrule
Answer samples
&
Answer-state uncertainty
&
Do repeated generations disagree, or has the answer surface become
artificially stable? \\
\addlinespace[0.25em]
Retrieved passages
&
Evidence-state uncertainty
&
Does the retrieved text locally entail, contradict, or fail to support the
generated answer? \\
\addlinespace[0.25em]
Graph trace
&
Retrieval-state uncertainty
&
Which anchors, typed triples, and paths made the answer reachable, and where
does the graph abstain? \\
\bottomrule
\end{tabular}

\vspace{0.45em}
\begin{minipage}{0.96\linewidth}
\footnotesize
\textit{The families are not independent causes of correctness.}
They are separately logged views of the same KG-RAG run: generated text,
retrieved evidence, and symbolic retrieval state.
\end{minipage}

\caption{%
  \textbf{Uncertainty is treated as an audit problem over pipeline objects.}
  Answer-state metrics see only sampled answers.
  Evidence-state metrics compare the answer with retrieved passages.
  Retrieval-state diagnostics inspect the graph trace itself: anchors,
  subject--relation--object triples, paths, and abstentions.
  This is why the families are reported separately.
  This table defines the objects; Table~\ref{tab:failure_modes} states how
  each family behaves across retrieval regimes, and
  Table~\ref{tab:failure_taxonomy} maps their joint outcomes into a
  $2{\times}2$ failure taxonomy.%
}
\label{tab:decomposition}
\end{table}
